\documentclass[11pt]{article}

../macros/macros.tex

% for sums_recursion_trees.tex
% \usepackage{wrapfig}
% \usetikzlibrary{matrix,arrows}

% \usepackage{ifthen}

% \usepackage{pdfx}

% \newcommand{\first}[1]{\ensuremath{\operatorname{\sf e}({#1})}}
% \newcommand{\remark}[1]{\marginpar{\raggedright\small\emph{#1}}}

\begin{document}

% \setcounter{section}{1}

\doTitle{5}{Greedy Algorithms}{Greedy Algorithms}
\newpage


\section{Introduction}\label{greedy: sec: intro}\doHeader{Introduction}

An algorithm is \emph{greedy} if it builds a solution piece by piece.
Once all the pieces have been added, and the solution is complete, the algorithm returns the solution.
A guiding principle is that the algorithm must maintain the following ``greedy'' invariant at all times:

\medskip
\centerline
{\textsf{greedy invariant}: \em The current (partial) solution can be extended to a correct (full) solution.}
\medskip

To design the algorithm, a good strategy is to start by finding a rule for the \emph{first} greedy step,
and a proof that the first piece chosen by the rule is part of some correct solution.
Doing this helps build intuition.
Furthermore, in many cases, after the first step, the problem that remains to be solved
is a smaller instance of the same problem,
so the algorithm can simply recurse to solve that problem.
Or, if that doesn't work, we can take a different approach:
try to generalize the rule for the first step,
to find a rule for extending the partial solution by one piece at every step,
and a proof that the generalized rule maintains the greedy invariant.

We will also see an example that doesn't fit either of these patterns.

\inputSection{making_change}
\newpage

\inputSection{activity_selection}
\newpage

\inputSection{covering_by_intervals}
\newpage

\inputSection{Huffman}
\newpage

\inputSection{interval_partitioning}
\newpage

 \inputSection{minimum_wait_time}  % TODO: CHECK FOR CORRECTNESS
\newpage

\inputSection{Dijkstra}
\newpage

\inputSection{MST}
\newpage

\inputSection{90_exercises}

\codeSection{
  greedy_algorithms/library.py,
  greedy_algorithms/activity_selection.py,
  greedy_algorithms/Huffman.py,
  greedy_algorithms/Dijkstra.py,
  greedy_algorithms/min_spanning_tree.py}

\end{document}


% We use the following standard terms.
% Given an instance $I$,
% define a set $S$ to be a \emph{feasible solution (for $I$)}
% if it meets the constraint of being a pairwise-disjoint subset of $I$.
% A feasible solution is a correct solution if it has maximum size.
% Let $\opt I $ be the maximum size of any feasible solution:
% \[\opt I  ~=~\max\{ |S| : S\subseteq I \emph{ and $I$ is pairwise disjoint}\}
%   ~=~\max\{ |S| : S \emph{ is feasible for } I\}.\]
% Given any feasible solution $S$,
% say that $S$ is \emph{optimal (for $I$)} if $|S|=\opt I $.
% So, given $I$, the problem is to find a feasible solution of maximum size,
% that is, an optimal solution.


% To prove that a greedy algorithm is correct, we will always
% prove that \emph{the algorithm never makes a mistake}.
% That is, it never adds a piece that makes it impossible to reach a correct solution.
% In other words, we always prove that the algorithm maintains an invariant of the following form:

% \medskip 
% \centerline{\textsf{greedy invariant}:
%   \emph{The current partial solution can be extended (somehow) to a full, correct solution.}}
% \medskip 

%%% Local Variables:
%%% mode: latex
%%% TeX-master: t
%%% End:
